\section{Introdução}

A pesquisa examina a construção, a apresentação e a circulação de indicadores de violência letal no Brasil, considerando os anos-base de 2016 a 2024. Seu ponto de partida é a distinção entre registros administrativos, produtos que os consolidam e materiais que comunicam seus resultados. A comparação de entidades por uma métrica comum envolve escolhas sociais de equivalência, como discutem \citeonline{espeland1998comensuracao}. No projeto, essa perspectiva orienta a reconstrução da proveniência dos indicadores, sem presumir que fontes com categorias próximas sejam automaticamente intercambiáveis. Este documento constitui um esqueleto didático, sem resultados empíricos concluídos.

A pergunta que orienta o estudo é: de que forma escolhas de construção e apresentação de indicadores de violência letal se relacionam com sua visibilidade institucional e midiática e com sua mobilização em processos federais de atenção e priorização estatal, considerando os anos-base de 2016 a 2024? O objetivo geral é analisar essas relações por meio de pesquisa documental, longitudinal e de abordagem mista. O percurso distingue produção estatística, comunicação institucional, cobertura jornalística e mobilização estatal federal. A investigação não medirá percepção individual da população nem converterá correspondências documentais ou temporais em demonstração de causalidade.
